\begingroup
\begin{center}\textbf{Abstract}\end{center}
\fontsize{9.1}{10.55}\selectfont
\setlength{\parskip}{1.2pt}
\setlength{\abovedisplayskip}{3pt}
\setlength{\belowdisplayskip}{3pt}
The Continuum Hypothesis is proved in the genealogical closure constructed
by Triadic Modular Holography (HMT): every internal subset of its real line is
countable or has the cardinality of the complete line. HMT is the formal core
composing three objects. All Possible Paths (APP) organizes sum and product
evaluations on a modular grid whose Cayley graph is the skeleton of a
two-dimensional toroidal cell complex. TRIT is the oriented ternary signature
\(\{-1,0,+1\}\). The Topological Phase Kernel (TPK) selects and transports
transitions, preserves their memory, and prolongs compatible histories. The
inverse limit \(X_\infty^{\mathrm{enr}}\) preserves value, sheet,
orientation, residue, quotient, carry, boundary, route, memory, and survival.
The joint discrete structure of the continuum and its Archimedean real line
are obtained as outputs of this action.

Let \(h\) be the effective code of the graphs of the APP, TRIT, and TPK
operators, their prolongations, their restrictions, and the two projections.
For each realized history \(m_\infty\), the genealogical closure
\(\mathfrak G_{\mathrm{HMT}}(h,m_\infty)\) is obtained by transfinite
iteration of finite definitions with parameters from earlier levels. The main
theorem establishes
\[
 \boxed{\mathfrak G_{\mathrm{HMT}}(h,m_\infty)
 \models 2^{\aleph_0}=\aleph_1}
\]
and, for every
\(A\in\mathcal P^{\mathfrak G}
(\mathbb R_{\mathrm{HMT}}^{\mathfrak G})\), proves the dichotomy
\[
 |A|\leq\aleph_0
 \qquad\text{or}\qquad
 |A|=|\mathbb R_{\mathrm{HMT}}^{\mathfrak G}|.
\]
Consequently, there is no intermediate cardinal between the natural numbers
and the real line within the integral continuum generated by HMT.

The proof links four steps. Exhaustive cylinder coverage establishes the
equinumerosity of the real line obtained through cylinder reduction with
\(\mathcal P^{\mathfrak G}(\omega)\). Genealogical
condensation constructs an injection
\(j_{\mathrm{HMT}}:\mathbb R_{\mathrm{HMT}}^{\mathfrak G}
\hookrightarrow\omega_1^{\mathfrak G}\). The coding of countable well-orders
and its ternary lift construct the injection in the opposite direction
\(k_{\mathrm{HMT}}\). Cantor--Bernstein yields cardinal equality, and the
cofinality argument extends the conclusion to every subset of the complete
closure. None of these steps presupposes the equality it proves.

The local fiber of \(3^6=729\) refinement subcylinders does not bound the
result: the coinductive law acts at every prescribed depth. The nonadic return
recovers the phase while carry and memory advance; it does not reset the
state. The same genealogy produces, as correlated results, the holonomy
\(\Gamma_9\), the integral dodecaphase register \(K\), the fourfold relation
\((\pi,\varphi,e,\alpha)\), and the
Paley--Witt--Golay--Leech--\(V^\natural\)--Monster incidence realization.
These outputs do not prove cardinal equality by numerical coincidence; they
certify the unity of the state whose Archimedean projection determines value
and whose incidence projection preserves structure.

Finally,
\(\mathfrak G_{\mathrm{HMT}}(h,m_\infty)=L[h,m_\infty]\) is proved as a
subsequent recognition of the hierarchy already constructed. It does not
generate the real line, does not select the injections, and does not enter
the cylinder reduction.

\medskip

% BEGIN R27_FRAMING_SUMMARY_EN
The description of the emitter is completed by the proof of its decimal
alphabet of \(42\) blocks, an explicit cycle of incidences between
regional half-emissions, and a contrast of ordered memory at fixed
histogram. These consequences specify the information preserved by the
finite readings preceding prolongation. Cardinal comparison retains
the domain and proofs proper to genealogical closure.
% END R27_FRAMING_SUMMARY_EN

\noindent\textbf{Keywords:} Triadic Modular Holography; APP; TRIT; TPK;
enriched state; discrete structure of the continuum; genealogical closure;
Continuum Hypothesis; cardinality; ordinality; coinduction.
\endgroup
